\chapter{标准函数}\label{ch:b1:functions}

分析学的用处，取决于人们所掌握的函数储备。
对于中学所学的这批函数——幂、指数、
对数、三角函数——本章再增加它们的反函数
（$\arcsin$、$\arccos$、$\arctan$）以及双曲函数族。
导数在高中阶段已被自由使用；
反函数的理论，
在\cref{ch:b1:continuity,ch:b1:derivative}中补全。

\section{指数、对数与幂}

\begin{proposition}[回顾与刻画]\label{prop:b1:functions:expln}
$\exp \colon \R \to \intoo{0}{+\infty}$ 与 $\ln \colon
\intoo{0}{+\infty} \to \R$ 为互逆双射，
严格递增，并有：
\[
\exp(x + y) = \exp x \exp y, \qquad \ln(xy) = \ln x + \ln y,
\qquad \exp' = \exp, \qquad \ln'(x) = \frac 1x .
\]
此外，$\exp$ 是\emph{唯一}的可导函数 $f \colon
\R \to \R$，满足 $f' = f$ 与 $f(0) = 1$。
\end{proposition}

\begin{proof}
回顾部分属于高中内容。关于唯一性，设 $f' = f$，
$f(0) = 1$，并令 $g(x) = f(x)\,\eu^{-x}$。则 $g' = f'\eu^{-x} -
f\eu^{-x} = 0$，故 $g$ 为常数，其值为 $g(0) = 1$，从而 $f = \exp$。
\end{proof}

\begin{definition}[一般幂]\label{def:b1:functions:powers}
对 $x > 0$ 及 $\alpha \in \R$，
$\;x^\alpha = \eu^{\alpha \ln x}$\index{幂函数}。对 $a > 0$，
$a \neq 1$，则 \emph{以 $a$ 为底的对数}为 $\log_a x =
\frac{\ln x}{\ln a}$，即 $x \mapsto a^x$ 的反函数。
\end{definition}

\begin{example}[解指数方程]\label{ex:b1:functions:expequations}
解 $2^x = 5^{\,x-1}$，求 $\R$ 中的解。两边均为正，故取
对数 --- 这是一步可逆的变形：
\[
x\ln2 = (x - 1)\ln5
\iff x(\ln2 - \ln5) = -\ln5
\iff x = \frac{\ln5}{\ln5 - \ln2} = \frac{\ln 5}{\ln\frac52}
\approx 1.756 .
\]
接着，解 $x^{\sqrt2} = 3$（$x > 0$），两边
取 $\frac1{\sqrt2}$ 次幂（即作用互逆双射）：
$x = 3^{1/\sqrt2} = \eu^{(\ln 3)/\sqrt2} \approx 2.175$。要点在于：
凡含幂的方程，都要经由 $\ln$ 与
$\exp$ 化开，因为定义 $x^\alpha = \eu^{\alpha\ln x}$
把一切幂的运算都化为指数的算术 ---
但只在该定义适用的 $x > 0$ 范围内成立。
\end{example}

\begin{example}[倍增时间]\label{ex:b1:functions:doubling}
某量每步增长 $3\%$，经过 $n$ 步后，它变为原来的
$(1.03)^n$ 倍。何时翻倍？解 $(1.03)^n
\geq 2$：
\[
n\ln(1.03) \geq \ln2
\iff n \geq \frac{\ln 2}{\ln 1.03}
= \frac{0.6931}{0.02956} \approx 23.4 ,
\]
故第一次翻倍发生在第 $24$ 步。（金融界所谓的
“$72$ 法则”，即以
$72$ 除以百分数表示的增长率来估计倍增时间，这正是
对 $\ln(1 + x) \approx x$ 作近似后的计算，其误差见
\cref{ch:b1:taylor}。）
指数过程最好借助其对数来思考：在
该尺度上，增长是线性的，问题便化为除法。
\end{example}

\begin{example}[$2^{2026}$ 有几位？]\label{ex:b1:functions:digits}
整数 $N \geq 1$ 的十进制位数为
$\floor{\log_{10} N} + 1$（事实上，$N$ 恰有 $d$ 位，当且仅当
$10^{d-1} \leq N < 10^d$，即 $d - 1 \leq \log_{10}N < d$）。对
$N = 2^{2026}$：
\[
\log_{10} 2^{2026} = 2026\,\log_{10}2
= 2026 \times 0.301030 = 609.887 ,
\]
故 $2^{2026}$ 有 $610$ 位。小数部分还带来
一份额外收获：$10^{0.887} \approx 7.7$，所以这个数的\emph{首位}为
$7$。一次乘法便回答了一个关于没人会真正写出的数
的问题 --- 对数把乘法意义上的大小压缩为
加法意义上的大小，这正是它们全部的历史用意
（\cref{ex:b1:structures:expmorphism}使这句话得到精确表述）。
\end{example}

\begin{proposition}[幂的运算法则与增长比较]\label{prop:b1:functions:powerrules}
对 $x, y > 0$ 及 $\alpha, \beta \in \R$，有：
\[
x^{\alpha+\beta} = x^\alpha x^\beta, \quad
(x^\alpha)^\beta = x^{\alpha\beta}, \quad
(xy)^\alpha = x^\alpha y^\alpha, \quad
(x^\alpha)' = \alpha\, x^{\alpha - 1}.
\]
增长阶\index{增长比较}，当 $x \to +\infty$ 时，对任意
$\alpha > 0$ 与 $\beta > 0$，有：
\[
\frac{(\ln x)^\beta}{x^\alpha} \longrightarrow 0,
\qquad
\frac{x^\alpha}{\eu^{\beta x}} \longrightarrow 0 .
\]
\end{proposition}

\begin{proof}
恒等式由 $\exp$ 与 $\ln$ 的相应恒等式经
定义转写而来。具体看第二式（最不显然的一式）：$x^\alpha
> 0$，且 $\ln(x^\alpha) = \alpha\ln x$（对定义取
$\ln$ 即得），于是
\[
(x^\alpha)^\beta = \eu^{\beta\ln(x^\alpha)}
= \eu^{\beta\alpha\ln x} = x^{\alpha\beta} ;
\]
第一式与第三式则是同样的两行展开，借助
$\exp(u + v) = \exp u\exp v$ 与 $\ln(xy) = \ln x + \ln y$ 即可。
求导由链式法则给出：$(\eu^{\alpha\ln x})' =
\frac{\alpha}{x} \eu^{\alpha\ln x}$。

比较：由 $\frac{\ln t}{t} \to 0$（在高中
分册中已证明），取 $t = x^{\alpha/\beta}$：
得 $\frac{\ln x}{x^{\alpha/\beta}} \to 0$，再作
$\beta$ 次幂。对第二式，在第一式中取 $x = \ln u$。
\end{proof}

\begin{example}[每位读者都应掌握的两个极限]\label{ex:b1:functions:xx}
求 $\lim_{x \to 0^+} x^x$ 与 $\lim_{x \to +\infty}
x^{1/x}$ 的值？两个幂都借助指数函数定义，因此
应先处理指数：
\[
x^x = \eu^{x\ln x}, \qquad
x\ln x = -\frac{\ln(1/x)}{1/x} \longrightarrow 0
\quad (x \to 0^+),
\qquad\text{从而 } x^x \longrightarrow \eu^0 = 1 ;
\]
且 $x^{1/x} = \eu^{(\ln x)/x} \to \eu^0 = 1$，当 $x \to +\infty$
时，直接由增长比较可得。要点在于：形如 $0^0$ 或 $\infty^0$ 的不定型幂，
\emph{总是}借助
改写 $u^v = \eu^{v\ln u}$ 并分析乘积 $v \ln u$，
--- 绝不靠分别观察底数与指数去猜测。
决定两个极限的函数 $x \mapsto \frac{\ln x}{x}$
在本章的周末问题中被详尽研究。
\end{example}

\begin{omfigure}
\begin{tikzpicture}
\begin{axis}[omaxis, width=9.5cm, height=6cm,
    xmin=0, xmax=5.2, ymin=-1.8, ymax=8.5,
    xtick={1,2,3,4,5}, ytick={2,4,6,8}]
  \addplot[omMeth, thick, domain=0.02:5, samples=200] {ln(x)}
    node[pos=0.97, below, font=\small] {$\ln x$};
  \addplot[omDef, thick, domain=0:5, samples=100] {sqrt(x)}
    node[pos=0.9, above, font=\small] {$\sqrt x$};
  \addplot[omProp, thick, domain=0:5] {x}
    node[pos=0.72, left, font=\small] {$x$};
  \addplot[omThm, thick, domain=0:2.14, samples=120] {exp(x)}
    node[pos=0.85, left, font=\small] {$\eu^x$};
\end{axis}
\end{tikzpicture}

{\small 下图示出
\cref{prop:b1:functions:powerrules} 所述的增长尺度：靠近右边缘时，
$\ln x$ 仅刚过 $1.6$，而 $\eu^x$ 已移出画幅。每个比值
（对数比幂、幂比指数）都趋于 $0$ ---
图形只是暗示，证明中的代换才使之
精确。}
\end{omfigure}

\begin{remark}[幂与对数的常见误区]\label{rem:b1:functions:pitfalls}
\begin{enumerate}
  \item \emph{定义域。} 无理指数 $x^\alpha$（$\alpha$ 为无理数）要求
        $x > 0$；$(-8)^{1/3}$ 最好避免，而改用“
        $-8$ 的实立方根”，因为规则
        $(x^\alpha)^\beta = x^{\alpha\beta}$ 在负数上悄然失效，
        例如 $\bigl((-8)^2\bigr)^{1/6} = 2 \neq -2$。
  \item \emph{$\sqrt{x^2} = \abs x$}，而非 $x$：忘记
        取绝对值是丢失负数解的
        经典来源。
  \item \emph{$\ln(xy) = \ln x + \ln y$ 要求 $x, y > 0$。}对于
        负的自变量，$\ln(xy)$ 可以有定义，而
        右端却无定义。
  \item \emph{哪个函数增长？} 在 $x^\alpha$ 中，
        \emph{底数}变化；在 $a^x$ 中则是\emph{指数}变化。
        二者的增长速度截然不同
        （\cref{prop:b1:functions:powerrules}），而混合型
        表达式如 $x^{\ln x}$ 或 $x^{1/x}$，
        必须先改写为 $\eu^{(\cdot)}$ 再作任何推理 ---
        如\cref{ex:b1:functions:xx}所示。
\end{enumerate}
\end{remark}

\section{反三角函数}

\begin{definition}[$\arcsin$、$\arccos$、$\arctan$]\label{def:b1:functions:arc}
下列限制
\[
\sin \colon \intcc{-\tfrac\pi2}{\tfrac\pi2} \to \intcc{-1}{1},
\qquad
\cos \colon \intcc{0}{\pi} \to \intcc{-1}{1},
\qquad
\tan \colon \intoo{-\tfrac\pi2}{\tfrac\pi2} \to \R
\]
都是严格单调的双射。其逆映射分别记作
$\arcsin$\index{反正弦}、$\arccos$\index{反余弦}与
$\arctan$\index{反正切}。例如，$y = \arcsin x$
（$x \in \intcc{-1}{1}$）是\emph{唯一的}角，它位于
$\intcc{-\frac\pi2}{\frac\pi2}$ 中，其正弦为 $x$。
\end{definition}

\begin{omfigure}
\begin{tikzpicture}
\begin{axis}[omaxis, width=5.6cm, height=5.6cm,
    xmin=-1.25, xmax=1.25, ymin=-1.8, ymax=1.8,
    xtick={-1,1}, ytick={-1.5708,1.5708},
    yticklabels={$-\frac{\pi}{2}$,$\frac{\pi}{2}$}]
  \addplot[omDef, very thick, domain=-1:1, samples=200]
    {rad(asin(x))};
\end{axis}
\begin{scope}[xshift=5.3cm]
\begin{axis}[omaxis, width=5.6cm, height=5.6cm,
    xmin=-1.25, xmax=1.25, ymin=-0.4, ymax=3.4,
    xtick={-1,1}, ytick={1.5708,3.1416},
    yticklabels={$\frac{\pi}{2}$,$\pi$}]
  \addplot[omThm, very thick, domain=-1:1, samples=200]
    {rad(acos(x))};
\end{axis}
\end{scope}
\begin{scope}[xshift=10.6cm]
\begin{axis}[omaxis, width=5.6cm, height=5.6cm,
    xmin=-4, xmax=4, ymin=-1.9, ymax=1.9,
    xtick={-2,2}, ytick={-1.5708,1.5708},
    yticklabels={$-\frac{\pi}{2}$,$\frac{\pi}{2}$}]
  \addplot[omMeth, very thick, domain=-4:4, samples=200]
    {rad(atan(x))};
  \addplot[dashed, gray, domain=-4:4] {1.5708};
  \addplot[dashed, gray, domain=-4:4] {-1.5708};
\end{axis}
\end{scope}
\end{tikzpicture}

{\small 从左到右依次为 $\arcsin$、$\arccos$、$\arctan$。每
一个函数的图像由限制后的直接函数经
关于直线 $y = x$ 的对称得到。}
\end{omfigure}

\begin{proposition}[导数]\label{prop:b1:functions:arcderiv}
在各自定义域的内部：
\[
\arcsin' x = \frac{1}{\sqrt{1 - x^2}},
\qquad
\arccos' x = \frac{-1}{\sqrt{1 - x^2}},
\qquad
\arctan' x = \frac{1}{1 + x^2}.
\]
\end{proposition}

\begin{proof}
下面用到的反函数求导法则将
在\cref{ch:b1:derivative}中证明：若 $f$ 为可导双射，且
$f' \neq 0$，则 $(f^{-1})'(x) = \frac{1}{f'(f^{-1}(x))}$。对
$\arcsin$，令 $y = \arcsin x$，
\[
\arcsin' x = \frac{1}{\cos y}
= \frac{1}{\sqrt{1 - \sin^2 y}} = \frac{1}{\sqrt{1 - x^2}},
\]
其中 $\cos y = +\sqrt{1 - \sin^2 y}$，因为 $y \in
\intoo{-\frac\pi2}{\frac\pi2}$ 蕴含 $\cos y > 0$。同理
$\arccos' x = \frac{1}{-\sin y}$，其中 $\sin y > 0$ 在
$\intoo{0}{\pi}$ 上成立，且 $\arctan' x = \frac{1}{1 + \tan^2 y} =
\frac{1}{1 + x^2}$，其中利用了 $\tan' = 1 + \tan^2$。
\end{proof}

\begin{example}[识别隐藏的常数]\label{ex:b1:functions:hiddenconst}
研究 $g(x) = \arctan\dfrac{1 - x}{1 + x}$ 在 $\intoo{-1}
{+\infty}$ 上。利用链式法则并作简单计算：
\[
g'(x) = \frac{1}{1 + \bigl(\frac{1-x}{1+x}\bigr)^2}\cdot
\frac{-(1+x) - (1-x)}{(1+x)^2}
= \frac{-2}{(1+x)^2 + (1-x)^2}
= \frac{-1}{1 + x^2} ,
\]
由于 $(1+x)^2 + (1-x)^2 = 2 + 2x^2$，故 $g' = -\arctan'$：
函数 $g + \arctan$ 在\emph{区间}
$\intoo{-1}{+\infty}$ 上导数恒为零，因而在其上为常数；其在 $x =
0$ 处的值为 $\arctan 1 + \arctan 0 = \frac\pi4$。结论：
\[
\arctan\frac{1 - x}{1 + x} = \frac\pi4 - \arctan x
\qquad (x > -1) .
\]
在 $\intoo{-\infty}{-1}$ 上，同样的求导计算仍成立，但
常数不同（$-\frac{3\pi}4$；需要
在 $x \to -\infty$ 时求极限）。要点在于：“导数为零蕴含
常数”是逐区间成立的论断 --- 正是
\cref{exo:b1:functions:6} 与
\cref{exo:b1:functions:10} 所依赖的细微之处。
\end{example}

\begin{proposition}[常用恒等式]\label{prop:b1:functions:arcidentities}
\begin{enumerate}
  \item 对 $x \in \intcc{-1}{1}$，$\arcsin x + \arccos x =
        \dfrac{\pi}{2}$。
  \item 对 $x > 0$，$\arctan x + \arctan\dfrac 1x = \dfrac{\pi}{2}$
        （$-\frac\pi2$，适用于 $x < 0$）。
  \item $\sin(\arccos x) = \cos(\arcsin x) = \sqrt{1 - x^2}$；
        $\;\tan(\arcsin x) = \dfrac{x}{\sqrt{1 - x^2}}$ 在
        $\abs x < 1$ 时成立。
\end{enumerate}
\end{proposition}

\begin{proof}
(1) 函数 $x \mapsto \arcsin x + \arccos x$ 的导数在
$\intoo{-1}{1}$ 上为零（由\cref{prop:b1:functions:arcderiv}可知），故该函数在
该区间上为常数，其值为 $\frac\pi2$，即该函数在 $0$ 处的值；端点
处的值 $\pm 1$ 可直接验证（$\frac\pi2 + 0$ 与 $-\frac\pi2 +
\pi$）。

(2) 同样的方法用于 $\intoo{0}{+\infty}$：导数为
$\frac{1}{1+x^2} + \frac{-1/x^2}{1 + 1/x^2} = \frac{1}{1+x^2} -
\frac{1}{x^2 + 1} = 0$，且 $x = 1$ 时，和为 $2\arctan 1 =
\frac\pi2$。当 $x < 0$ 时，利用 $\arctan$ 为奇函数。

(3) 令 $y = \arccos x \in \intcc{0}{\pi}$，则 $\sin y \geq 0$，故
$\sin y = \sqrt{1 - \cos^2 y} = \sqrt{1 - x^2}$。
对 $\cos(\arcsin x)$ 同理，正切公式则为两者之商。
\end{proof}

\begin{remark}
$\arcsin(\sin\theta) = \theta$ 成立，\emph{仅}当 $\theta \in
\intcc{-\frac\pi2}{\frac\pi2}$ 时成立，例如
$\arcsin(\sin\pi) = 0 \neq \pi$。另一方向的复合
$\sin(\arcsin x) = x$ 则在整个
$\intcc{-1}{1}$ 上成立。
\end{remark}

\begin{example}[回到主值区间]
\label{ex:b1:functions:wrap}
计算
\[
\arctan\Bigl(\tan\frac{3\pi}4\Bigr)
\qquad\text{和}\qquad
\arctan\Bigl(\tan\frac{17\pi}5\Bigr).
\]
做法：将所给的角换成
$\intoo{-\frac\pi2}{\frac\pi2}$ 中与之具有相同正切的唯一角，
即减去适当的 $\pi$ 的整数倍
（$\tan$ 的周期）。第一个：$\frac{3\pi}4 - \pi =
-\frac\pi4$，故答案为 $-\frac\pi4$。第二个：$\frac{17\pi}
5 - 3\pi = \frac{2\pi}5 \in \intoo{-\frac\pi2}{\frac\pi2}$，故
答案为 $\frac{2\pi}5$。该计算实际上是
以 $\pi$ 为除数对角度作欧几里得除法的伪装 --- 而
对 $\arcsin$ 的类似做法（反射至
$\intcc{-\frac\pi2}{\frac\pi2}$，周期 $2\pi$）与 $\arccos$
的做法（反射至 $\intcc0\pi$）则分别用于 \cref{exo:b1:functions:1} 与
\cref{exo:b1:functions:10} 中的分段答案。
\end{example}

\begin{example}[安全地求反正切之和]\label{ex:b1:functions:arctansum}
下面证明
\[
\arctan\frac12 + \arctan\frac15 + \arctan\frac18 = \frac\pi4 .
\]
这里要用两个要点：正切的和角公式，以及 --- 初学者常忘的一步 ---
对和的\emph{定位}。首先，
$\tan(u + v) = \frac{\tan u + \tan v}{1 - \tan u\tan v}$，
取 $u = \arctan\frac12$，$v = \arctan\frac15$，可得
\[
\tan(u + v) = \frac{\frac12 + \frac15}{1 - \frac1{10}}
= \frac{7/10}{9/10} = \frac79,
\qquad\text{则}\qquad
\tan\Bigl(u + v + \arctan\frac18\Bigr)
= \frac{\frac79 + \frac18}{1 - \frac7{72}}
= \frac{65/72}{65/72} = 1 .
\]
其次，三个角中的每一个都位于 $\intoo0{\frac\pi4}$ 内
（它们的自变量都小于 $1$），故和位于
$\intoo0{\frac{3\pi}4}$ 内；正切为 $1$ 的角只有
$\frac\pi4$。若不做定位，结论只能是
“$\frac\pi4$ 加上 $\pi$ 的整数倍” --- 这只是半个证明。同样的
两步式思路也贯穿\cref{exo:b1:functions:8}与
\cref{exo:b1:functions:12}。
\end{example}

\begin{example}[一个反正切方程]\label{ex:b1:functions:arctaneq}
解方程 $\arctan x + \arctan 2x = \dfrac\pi4$。先作定位：左边
的符号与 $x$ 相同（两项都是），因此任一解满足
$x > 0$，于是其和位于 $\intoo0\pi$ 内。取正切
（在任何长为 $\pi$ 的区间上都不是单射，但结合
定位便已足够）：由和角公式得
\[
\tan\bigl(\arctan x + \arctan 2x\bigr)
= \frac{3x}{1 - 2x^2} = 1
\iff 2x^2 + 3x - 1 = 0
\iff x = \frac{-3 \pm \sqrt{17}}4 .
\]
负根由定位舍去；对于
正的候选值 $x_0 = \frac{\sqrt{17} - 3}4 \approx 0.28$，
其和位于 $\intoo0\pi$ 内且正切为 $1$，满足此条件的角只有
$\frac\pi4$（在 $\intoo{\frac\pi2}\pi$ 上，正切
为负）：$x_0$ 是唯一解。注意
论证的结构：取正切可能引入新解，但绝不会丢失解，
因此应先解多项式方程，\emph{再}用
定位筛选——这与两边平方解方程时相同的
逐一检验原则。
\end{example}

\section{双曲函数}

\begin{definition}[$\cosh$, $\sinh$, $\tanh$]\label{def:b1:functions:hyperbolic}
对于 $x \in \R$：
\[
\cosh x = \frac{\eu^x + \eu^{-x}}{2},
\qquad
\sinh x = \frac{\eu^x - \eu^{-x}}{2},
\qquad
\tanh x = \frac{\sinh x}{\cosh x}
\]
（\emph{双曲余弦、双曲正弦、双曲正切}\index{双曲
函数}）。$\cosh$ 为偶函数，$\sinh$ 与 $\tanh$ 为奇函数。
\end{definition}

\begin{proposition}[基本性质]\label{prop:b1:functions:hyprules}
对所有 $x, y \in \R$：
\begin{enumerate}
  \item $\cosh^2 x - \sinh^2 x = 1$;
  \item $\cosh' = \sinh$, $\;\sinh' = \cosh$, $\;\tanh' = 1 - \tanh^2
        = \dfrac{1}{\cosh^2}$;
  \item $\cosh(x+y) = \cosh x\cosh y + \sinh x \sinh y$ 与
        $\sinh(x+y) = \sinh x \cosh y + \cosh x \sinh y$;
  \item $\sinh$ 是 $\R$ 到 $\R$ 上的严格递增双射；
        $\cosh$ 限制在 $\R_+$ 上后，是严格递增的
        到 $\intco{1}{+\infty}$ 上的双射；$\tanh$ 是严格
        递增的双射，将 $\R$ 映到 $\intoo{-1}{1}$ 上。
\end{enumerate}
\end{proposition}

\begin{proof}
(1) $\bigl(\frac{\eu^x + \eu^{-x}}{2}\bigr)^2 - \bigl(\frac{\eu^x -
\eu^{-x}}{2}\bigr)^2 = \frac{(\eu^{2x} + 2 + \eu^{-2x}) - (\eu^{2x} -
2 + \eu^{-2x})}{4} = 1$.

(2) 对定义式求导；对于 $\tanh$，由商
法则得 $\frac{\cosh^2 - \sinh^2}{\cosh^2}$，它既
等于 $1 - \tanh^2$，又等于 $\frac{1}{\cosh^2}$（由 (1)）。

(3) 利用定义展开右端。具体地，
对第一式：
\[
\cosh x\cosh y + \sinh x\sinh y
= \frac{(\eu^x + \eu^{-x})(\eu^y + \eu^{-y})
+ (\eu^x - \eu^{-x})(\eu^y - \eu^{-y})}{4} ;
\]
在八个乘积中，四个“混合”项
（$\eu^{x-y}$、$\eu^{y-x}$）成对相消，而
$\eu^{x+y}$ 与 $\eu^{-(x+y)}$ 各出现两次：
总和为 $\frac{2\eu^{x+y} + 2\eu^{-(x+y)}}{4} = \cosh(x+y)$。
正弦式的证明完全相同，只是改为保留
混合项。

(4) $\sinh' = \cosh \geq 1 > 0$，故 $\sinh$ 严格递增，
其极限为 $\pm\infty$（主导项为 $\pm\frac{\eu^{\abs
x}}{2}$），于是双射性由介值
定理推出（高中阶段已使用；系统论述见
\cref{ch:b1:continuity}）。在 $\R_+$ 上，$\cosh' = \sinh \geq 0$
仅在 $0$ 处为零：从 $\cosh 0 = 1$ 严格递增
到 $+\infty$。又 $\tanh' > 0$，其极限为 $\pm 1$（分别对应 $\pm\infty$）：
具体地写作：
\[
\tanh x = \frac{\eu^x - \eu^{-x}}{\eu^x + \eu^{-x}}
= \frac{1 - \eu^{-2x}}{1 + \eu^{-2x}}
\xrightarrow[x \to +\infty]{} 1 ,
\]
分子分母同除以 $\eu^x$ 即得，再由奇性
得极限 $-1$（位于 $-\infty$ 处）；于是严格递增的
函数 $\tanh$ 因此将 $\R$ 映到 $\intoo{-1}1$ 上。
\end{proof}

\begin{omfigure}
\begin{tikzpicture}
\begin{axis}[omaxis, width=6.4cm, height=6cm,
    xmin=-2.6, xmax=2.6, ymin=-3.6, ymax=3.9,
    xtick={-2,-1,1,2}, ytick={-3,-2,-1,1,2,3}]
  \addplot[omDef, very thick, domain=-2.05:2.05, samples=120]
    {cosh(x)} node[pos=0.1, right, font=\small] {$\cosh$};
  \addplot[omThm, very thick, domain=-2.05:2.05, samples=120]
    {sinh(x)} node[pos=0.93, above left, font=\small] {$\sinh$};
  \addplot[dashed, gray, domain=0:2.5] {exp(x)/2};
\end{axis}
\begin{scope}[xshift=7.4cm]
\begin{axis}[omaxis, width=6.4cm, height=6cm,
    xmin=-3.2, xmax=3.2, ymin=-1.5, ymax=1.5,
    xtick={-2,-1,1,2}, ytick={-1,1}]
  \addplot[omMeth, very thick, domain=-3:3, samples=150]
    {tanh(x)} node[pos=0.9, above left, font=\small] {$\tanh$};
  \addplot[dashed, gray, domain=-3:3] {1};
  \addplot[dashed, gray, domain=-3:3] {-1};
\end{axis}
\end{scope}
\end{tikzpicture}

{\small 左图：$\cosh$ 与 $\sinh$，渐近地贴合于
$\frac{\eu^x}{2}$（虚线）。右图：$\tanh$，自 $-1$ 递增至
$1$。}
\end{omfigure}

\begin{remark}[为什么叫“双曲”？]
点 $(\cosh t, \sinh t)$ 沿 $x > 0$ 的那支双曲线
$x^2 - y^2 = 1$ 运行（由
\cref{prop:b1:functions:hyprules}~(1)），正如 $(\cos t, \sin
t)$ 沿圆 $x^2 + y^2 = 1$ 运行。每个圆函数恒等式
都有对应的双曲形式，符号的变化由
$\sinh^2 \leftrightarrow -\sin^2$ 决定。
\end{remark}

\begin{example}[$\tanh$ 的加法公式]\label{ex:b1:functions:tanhadd}
将
\cref{prop:b1:functions:hyprules}~(3)中的两个加法公式都除以 $\cosh x\cosh y$：
\[
\tanh(x + y)
= \frac{\sinh x\cosh y + \cosh x\sinh y}
{\cosh x\cosh y + \sinh x\sinh y}
= \frac{\tanh x + \tanh y}{1 + \tanh x\,\tanh y} ,
\]
即正切加法公式的双曲类似式 --- 其中
$+$ 对应三角情形中的 $-$。附带的收获是：由 $\abs{\tanh}
< 1$，右端是一条“速度叠加”法则，
它永不越出 $\intoo{-1}1$：若 $u, v \in \intoo{-1}1$，则
$\frac{u + v}{1 + uv} \in \intoo{-1}1$ 亦然（记 $u = \tanh
a$，$v = \tanh b$，由双射性可作此设，并将公式
倒过来用）。直接用代数方法验证，不借助双曲函数，
是一道略显繁琐的练习；若以
$\tanh$ 作参数化，则只需一行 --- 这与圆函数
为单位圆提供的策略如出一辙。
\end{example}

\begin{proposition}[反双曲函数]\label{prop:b1:functions:invhyp}
由\cref{prop:b1:functions:hyprules}~(4)所保证的反函数
均有如下显式表达式：
\[
\operatorname{arsinh} x = \ln\bigl(x + \sqrt{x^2 + 1}\bigr)
\ (x \in \R),
\qquad
\operatorname{arcosh} x = \ln\bigl(x + \sqrt{x^2 - 1}\bigr)
\ (x \geq 1),
\]
\[
\operatorname{artanh} x = \frac 12 \ln\frac{1 + x}{1 - x}
\ (\abs x < 1),
\]
其导数依次为 $\frac{1}{\sqrt{x^2+1}}$、$\frac{1}{\sqrt{x^2-1}}$
（$x > 1$）与 $\frac{1}{1 - x^2}$。
\end{proposition}

\begin{proof}
对 $\operatorname{arsinh}$：解方程 $x = \sinh y = \frac{\eu^y -
\eu^{-y}}{2}$。令 $u = \eu^y > 0$，得 $u^2 - 2xu - 1 = 0$，故 $u = x
+ \sqrt{x^2 + 1}$（另一根 $x - \sqrt{x^2+1}$ 为负），且 $y =
\ln(x + \sqrt{x^2+1})$。另外两式的计算完全类似
（$u^2 - 2xu + 1 = 0$，对应 $\operatorname{arcosh}$，取
$\geq 1$ 的根）；对 $\operatorname{artanh}$ 只需化简两行即可。
求导：对对数式直接求导，例如
\[
\bigl(\ln(x + \sqrt{x^2+1})\bigr)'
= \frac{1 + \frac{x}{\sqrt{x^2+1}}}{x + \sqrt{x^2+1}}
= \frac{1}{\sqrt{x^2 + 1}} . \qedhere
\]
\end{proof}

\begin{example}[显式公式的应用]\label{ex:b1:functions:arcoshtwo}
由显式公式，方程 $\cosh t = 2$ 在 $t \geq 0$ 时的解为
$t = \operatorname{arcosh} 2 = \ln(2 + \sqrt3) \approx
1.317$；由偶性，另一解为 $-t$。事实上
$\ln(2 - \sqrt3) = \ln\frac1{2 + \sqrt3} = -\ln(2 + \sqrt3)$：
$u = \eu^t$ 满足二次方程 $u^2 - 4u + 1 = 0$，其两根
互为倒数（其积为 $1$，正是韦达定理所要求的）。这个小例子
体现了一般的规律：双曲方程
可化为关于 $\eu^t$ 的二次方程，而 $t \mapsto -t$ 的对称性
在二次方程中表现为 $u \mapsto 1/u$ 的对称性 ---
解 \cref{exo:b1:functions:7} 时值得牢记。
\end{example}

\begin{method}[选取恰当的原函数形式]\label{met:b1:functions:primitive}
求积分时需熟记的三个基本形式
（在\cref{ch:b1:integration}中）：
\[
\int \frac{\dd x}{1 + x^2} = \arctan x + C,
\qquad
\int \frac{\dd x}{\sqrt{1 - x^2}} = \arcsin x + C,
\qquad
\int \frac{\dd x}{\sqrt{x^2 + 1}} = \operatorname{arsinh} x + C,
\]
对一般的二次式，先配方，
再作伸缩，化归为上述形式。
\end{method}

\begin{remark}[本章的用途]\label{rem:b1:functions:whereused}
本章是此后全部分析的常用工具。
\cref{prop:b1:functions:powerrules}的增长尺度决定
全书中的收敛问题，
见\cref{ch:b1:seq,ch:b1:series}；反三角函数的导数公式
\cref{prop:b1:functions:arcderiv}与
\cref{prop:b1:functions:invhyp}，它们是
在\cref{ch:b1:integration}中最常需要的原函数，借助
\cref{met:b1:functions:primitive}求得；双曲函数
参数化 $y'' = y$ 的解，
在\cref{ch:b1:diffeq}中，其方式正如圆函数参数化 $y'' = -y$ 的解。
函数 $\exp$ 由 $f' = f$，
$f(0) = 1$ 的刻画（见\cref{prop:b1:functions:expln}）是
线性微分方程理论的一维种子，而
悬链线 $y = \cosh x$ 重新出现在
\cref{ch:b1:curves}的平面曲线之中。
\end{remark}

\begin{remark}[插曲：反函数程序]\label{rem:b1:functions:interlude}
本章用一个程序做了四次：限制函数，
直到它成为严格单调的双射，写出其反函数，
并将每个公式经它转换。每一步中所\emph{用到}的事实——
区间上的连续严格单调函数是到某区间的双射，
其反函数连续，且在临界点之外可导——
这些事实是凭高中阶段的信任借来的。
这笔欠账在本分册中偿还：
\cref{ch:b1:continuity}证明了双射结论
（单调反函数定理，依据介值定理），
而\cref{ch:b1:derivative}证明了导数法则
$(f^{-1})' = 1/(f' \circ f^{-1})$，它无声地推出了
\cref{prop:b1:functions:arcderiv}与
\cref{prop:b1:functions:invhyp}的全部公式。读这些章节时
带着 $\arcsin$ 与 $\operatorname{arsinh}$，
将它们视作该理论本为证明而建立的典型例子——
这正是绕开表面上循环论证的预期方式。
\end{remark}

\section{练习}

\begin{exercise}[$\star$]\label{exo:b1:functions:1}
不用计算器计算：$\arcsin\frac{\sqrt 3}{2}$；
$\;\arccos\bigl(-\frac 12\bigr)$; $\;\arctan(-1)$;
$\;\arcsin\bigl(\sin\frac{5\pi}{6}\bigr)$;
$\;\arccos\bigl(\cos\frac{7\pi}{4}\bigr)$.
\end{exercise}

\begin{exercise}[$\star$]\label{exo:b1:functions:2}
给出 $f(x) = \arcsin(2x - 1)$ 的定义域，并
在其定义域内求 $f'$。对 $g(x) = \arctan\sqrt{x}$ 提出同样的问题。
\end{exercise}

\begin{exercise}[$\star$]\label{exo:b1:functions:3}
证明：对任意 $x \in \R$，$\cosh x + \sinh x = \eu^x$，并且
$(\cosh x + \sinh x)^n = \cosh nx + \sinh nx$ 对每个 $n \in \Z$ 成立
（双曲形式的棣莫弗公式）。
\end{exercise}

\begin{exercise}[$\star$]\label{exo:b1:functions:4}
将 $\cos(2\arcsin x)$ 与 $\sin(2\arctan x)$ 化为
$x$ 的代数式。
\end{exercise}

\begin{exercise}[$\star$]\label{exo:b1:functions:5}
对充分大的 $x$，将下列函数
$x^{100}$, $\;\eu^{\sqrt x}$, $\;(\ln x)^{1000}$, $\;\eu^{x/100}$,
$\;x^{\ln x}$，从增长最慢到最快排列，并
依据\cref{prop:b1:functions:powerrules}说明理由。
\end{exercise}

\begin{exercise}[$\star\star$]\label{exo:b1:functions:6}
研究函数 $f(x) = \arctan\dfrac{2x}{1 - x^2}$ 在其定义域上的性质：
求 $f'$，与 $(2\arctan x)'$ 比较，并将 $f(x)$
在定义域的三个区间上分别用 $\arctan x$ 表示。
\end{exercise}

\begin{exercise}[$\star\star$]\label{exo:b1:functions:7}
在 $\R$ 中解方程 $\;\cosh x = 2$；再解 $5\cosh x - 4 \sinh x = 3$
（解用对数表示）。\emph{第二题提示：
把所有式子都用 $u = \eu^x$ 写出。}
\end{exercise}

\begin{exercise}[$\star\star$]\label{exo:b1:functions:8}
证明恒等式 $\arctan\frac{1}{2} + \arctan\frac{1}{3} =
\frac{\pi}{4}$，再证马青公式\index{马青公式}
\[
4\arctan\frac 15 - \arctan\frac{1}{239} = \frac{\pi}{4}.
\]
\emph{提示：用加法公式计算等式两边的正切，
并判断两边各自所在的区间。}
\end{exercise}

\begin{exercise}[$\star\star$]\label{exo:b1:functions:9}
证明：对任意 $x \geq 0$，$\;x - \dfrac{x^3}{6} \leq \sin x \leq
x$ \emph{（研究各差函数的逐阶导数）}，并由此
推出 $\lim_{x \to 0^+} \frac{x - \sin x}{x^3} = \frac 16$ 是
合理的——该极限本身将在
\cref{ch:b1:taylor}中确立。
\end{exercise}

\begin{exercise}[$\star\star\star$]\label{exo:b1:functions:10}
对 $x \in \intcc{-1}{1}$，令
$h(x) = \arcsin\bigl(2x\sqrt{1 - x^2}\,\bigr) - 2\arcsin x$。人们或许会猜想
$h = 0$，这由恒等式 $\sin 2y = 2\sin y\cos y$ 而来——
但 $h$ \emph{并非}恒等于零。
在其存在导数的开区间上求 $h'$，在适当选取的点处计算 $h$，
并给出 $h$ 在 $\intcc{-1}{1}$ 上的完整分段常值描述。
\end{exercise}

\begin{exercise}[$\star\star\star$]\label{exo:b1:functions:11}
（古德曼函数）令 $g(x) = \arctan(\sinh x)$，其中 $x \in \R$。证明
$g$ 是奇函数，且是由 $\R$ 到
$\intoo{-\frac\pi2}{\frac\pi2}$ 上的严格递增双射，$g'(x) = \frac{1}{\cosh x}$，
且 $\tan g(x) = \sinh x$，$\;\sin g(x) = \tanh x$，
$\;\cos g(x) = \frac{1}{\cosh x}$：函数 $g$ 将圆三角
与双曲三角联系起来，无需借助复数。
\end{exercise}

\begin{exercise}[$\star\star$]\label{exo:b1:functions:12}
证明：
\[
\arctan 1 + \arctan 2 + \arctan 3 = \pi .
\]
\emph{提示：先计算 $\arctan 2 + \arctan 3$，
并像 \cref{ex:b1:functions:arctansum} 那样确定其取值范围；注意此处
加法公式的分母为负。}
\end{exercise}

\section{问题：方程 $x^y = y^x$}

\begin{problem}\label{pb:b1:functions:1}
哪些正数对满足 $x^y = y^x$？人人都知道
一个看似偶然的例子 $2^4 = 4^2 = 16$；本题将说明
这一切并非偶然。整个方程受
单一函数的变化情况所支配：
\[
f(t) = \frac{\ln t}{t} \qquad (t > 0),
\]
研究该函数可得：方程的完整解集（一条对角线，加上
经过 $(\eu, \eu)$ 的一支曲线），该分支的有理参数化，
$(2, 4)$ 是其\emph{唯一}整数点这一事实，
以及全部有理点的分类 --- 此外，作为副产品还有
$\eu^\pi$ 与 $\pi^\eu$ 的比较，以及
$\bigl(1 + \frac1n\bigr)^n$ 单调收敛于 $\eu$。
\index{幂函数}

\medskip
\noindent\textbf{第一部分 --- 函数 $f(t) = \ln t / t$。}
\begin{enumerate}
  \item 说明 $f$ 在 $\intoo0{+\infty}$ 上可导，
        计算 $f'$，并列出变化表：$f$
        在 $\intoc0\eu$ 上严格递增，在
        $\intco\eu{+\infty}$ 上严格递减，最大值为 $f(\eu) = \frac1\eu$。
  \item 求 $f$ 在 $0^+$ 与 $+\infty$ 处的极限
        （见 \cref{prop:b1:functions:powerrules}），$f$ 的符号
        （在 $\intoo01$ 上为负，在 $1$ 处为零，之后为正），
        并描绘其图像。
  \item 直接计算验证 $f(2) = f(4)$。（请记住
        这一等式，整个问题都由它衍生。）
  \item 设 $c \in \R$，按 $c$ 的取值讨论
        方程 $f(t) = c$ 解的个数：当 $c \leq
        0$ 时恰有一解；两解的情形（一解在 $\intoo1\eu$ 内，另一解在
        $\intoo\eu{+\infty}$ 内）对应 $0 < c < \frac1\eu$；
        当 $c = \frac1\eu$ 时恰有一解；当 $c > \frac1\eu$ 时无解。
  \item 由此推出：$\eu^x > x^\eu$ 对一切 $x > 0$ 且 $x \neq
        \eu$ 成立；特别地，判定 $\eu^\pi$ 与
        $\pi^\eu$ 哪一个更大。
\end{enumerate}

\medskip
\noindent\textbf{第二部分 --- 方程及其曲线。}在本部分中，
$x, y > 0$。
\begin{enumerate}[resume]
  \item 证明 $x^y = y^x \iff f(x) = f(y)$。
  \item 由此推出解集的结构：全部对角点对 $(x, x)$；以及
        \emph{非平凡}点对（$x \neq
        y$），满足：两个坐标均 $> 1$，且其中一个
        位于 $\intoo1\eu$ 内，另一个位于
        $\intoo\eu{+\infty}$ 内。
  \item 对非平凡点对作参数化：记 $y = tx$，其中 $t
        > 0$，$t \neq 1$，证明 $x^y = y^x$ 蕴含
        \[
        x(t) = t^{\frac1{t-1}},
        \qquad
        y(t) = t\,x(t) = t^{\frac{t}{t-1}},
        \]
        反之，每个这样的点对都是解。
  \item 验证 $t = 2$ 给出 $(2, 4)$，并证明对称性
        $x(1/t) = y(t)$，$y(1/t) = x(t)$：取参数的倒数时
        两个坐标互换。
  \item 求 $x(t)$ 与 $y(t)$ 在 $t \to 1$ 时的极限
        （二者均趋于 $\eu$），在 $t \to +\infty$ 时（$x \to 1$，$y
        \to +\infty$），以及在 $t \to 0^+$ 时（$x \to +\infty$，$y \to
        1$）的极限。描述所得分支：一条渐近于直线
        $x = 1$ 与 $y = 1$ 的曲线，与对角线交于
        $(\eu, \eu)$。
  \item 证明 $t \mapsto x(t)$ 在
        $\intoo1{+\infty}$ 上严格递减。\emph{（研究 $h(t) = 1 - \frac1t - \ln
        t$，其符号决定 $\ln x(t) =
        \frac{\ln t}{t - 1}$ 的导数。）}
  \item 由此可知，非平凡解确定一个严格递减的双射
        $\varphi \colon \intoo1\eu \to
        \intoo\eu{+\infty}$ 满足 $x^{\varphi(x)} = \varphi(x)^x$，
        且 $\varphi$ 是该分支上的对合：
        $\varphi(\varphi(x)) = x$，只要两边有意义。
\end{enumerate}

\medskip
\noindent\textbf{第三部分 --- 整数点与有理点。}
\begin{enumerate}[resume]
  \item 证明 $(2, 4)$ 与 $(4, 2)$ 是 $x^y = y^x$ 仅有的非平凡
        \emph{整数}解。
  \item 对 $n \in \N^*$，取 $t = 1 +
        \frac1n$，代入参数化公式，证明
        \[
        x_n = \Bigl(1 + \frac1n\Bigr)^{n},
        \qquad
        y_n = \Bigl(1 + \frac1n\Bigr)^{n+1}
        \]
        它对每个 $n$ 都是非平凡的\emph{有理}解，
        且 $(x_1, y_1) = (2, 4)$。
  \item 反之，设 $(x, y)$ 为非平凡有理解，
        满足 $y > x$，记 $t = y/x = 1 + \frac rs$（既约形式，
        $r, s \in \N^*$）。利用 $x = t^{s/r}$，并承认
        质因数分解的唯一性（中学已熟悉；
        其证明见 \cref{ch:b1:arith}），证明若 $x$
        为有理数，则 $s$ 与 $s + r$ 都必须是
        整数的 $r$ 次幂。
  \item 用二项式定理证明 $a^r - b^r = r$ 没有
        满足 $a > b \geq 1$ 的整数解（$r \geq 2$），由此得出：
        $x^y = y^x$ 的有理解恰为第 14 题中的
        点对 $(x_n, y_n)$（及其
        互换）。
  \item 数值验证 $n = 2$ 的情形：计算 $f(9/4)$ 与
        $f(27/8)$，精确到小数点后四位，并检验二者相等。
\end{enumerate}

\medskip
\noindent\textbf{第 IV 部分 --- 收获。}
\begin{enumerate}[resume]
  \item 利用第 7 题与第 11 题，不作进一步计算，
        证明数列 $x_n = \bigl(1 +
        \frac1n\bigr)^n$ 严格递增，且 $x_n < \eu$，
        $y_n = \bigl(1 + \frac1n\bigr)^{n+1}$ 严格递减，且 $y_n > \eu$，
        并证明二者都收敛于
        $\eu$。（参数 $t_n = 1 + \frac1n$ 递减至
        $1$。）
  \item 建立 $1 < a < b$ 时的一般比较规律：
        若 $\eu \leq a$，则 $a^b > b^a$；若 $b \leq \eu$，则
        $a^b < b^a$；并通过两个例子 $(2, 3)$ 和
        $(2, 5)$ 表明，在混合情形 $a < \eu < b$ 下，
        两种结果确实都会出现。
  \item 设第 12 题中的对合 $\varphi$
        在 $\eu$ 处可导（事实上确实可导）。
        对恒等式 $\varphi(\varphi(x)) = x$ 在 $x = \eu$ 处求导，并推出
        $\varphi'(\eu) = -1$：该分支与对角线
        成直角相交。
  \item 求满足 $y = 3x$ 的唯一解对，写成闭形式，
        并借助 $f$ 数值检验，精确到小数点后四位。
  \item 在数 $\sqrt2$、$\sqrt[3]3$、$\sqrt[4]4$、
        $\sqrt[5]5$ 中，求出最大者，并指出哪两个数相等。
        \emph{（比较 $n^{1/n} =
        \eu^{f(n)}$。）}
\end{enumerate}

\medskip
\noindent\textbf{第 V 部分 --- 综合。}
\begin{enumerate}[resume]
  \item 描述方程 $x^y = y^x$ 在
        第一象限 $x, y > 0$ 内的完整解集——对角线与分支，
        它们的交点 $(\eu, \eu)$、渐近线、整数
        点 $(2, 4)$，以及在 $(\eu, \eu)$ 处聚集的有理点——
        写成一个能凭记忆画出草图的形式。
  \item 本题究竟在何处用到了：(i) 增长比较（见
        \cref{prop:b1:functions:powerrules}）；(ii)
        介值定理（借助双射的
        表述）；(iii) 在本阶段不加证明地承认的质因数分解唯一性？
        每项各用一句话回答。
  \item 寓意（用一小段文字）：一张单调性表
        解出了一个二元方程，并对其
        有理点作了分类，并证明了
        $\bigl(1 + \frac1n\bigr)^n$ 的单调收敛性。评述这种经济性，并指出
        各条线索在本卷后文中何处得到规模化应用，
        （数列见\cref{ch:b1:seq}，
        单调性表见\cref{ch:b1:derivative}，
        单调性表所缺的精度见\cref{ch:b1:taylor}）。
\end{enumerate}
\end{problem}
